\documentclass{harryopo-mathnotes}

\renewcommand{\mathtitle}{高等数学笔记}
\renewcommand{\mathauthor}{张三}
\date{2026 年 6 月 21 日}

\begin{document}

\newgeometry{top=3cm,bottom=2.5cm,left=4cm,right=4cm}
\maketitle
\thispagestyle{empty}
\cleardoublepage

\setcounter{tocdepth}{2}
\tableofcontents
\cleardoublepage

\strictpagecheck
\setcounter{page}{1}
\restoregeometry
\onehalfspacing

\section{函数与极限}

\subsection{函数的概念}

设 $x$ 和 $y$ 是两个变量，$D$ 是一个给定的数集。如果对于每个数 $x \in D$，变量 $y$ 按照一定的法则总有确定的数值和它对应，则称 $y$ 是 $x$ 的\textbf{函数}，记作 $y = f(x)$。

这里包含两个核心要素：

\begin{itemize}
  \item \textbf{定义域}：自变量 $x$ 的取值范围 $D$
  \item \textbf{对应法则}：$f$ 表示 $x$ 与 $y$ 之间的对应关系
\end{itemize}

\subsubsection{函数的表示法}

函数主要有三种表示法：

\begin{enumerate}
  \item \textbf{解析法}（公式法）：$y = x^2 + 2x + 1$
  \item \textbf{表格法}：用表格列出 $x$ 与 $y$ 的对应值
  \item \textbf{图像法}：用坐标系中的曲线表示函数关系
\end{enumerate}

\subsection{常见函数类型}

\begin{table}[htbp]
\centering
\begin{tabularx}{\textwidth}{>{\hsize=0.563\hsize\linewidth=\hsize}X>{\hsize=1.465\hsize\linewidth=\hsize}X>{\hsize=0.845\hsize\linewidth=\hsize}X>{\hsize=1.127\hsize\linewidth=\hsize}X}
\toprule
\textbf{类型} & \textbf{表达式} & \textbf{定义域} & \textbf{特点} \\
\midrule
---------- & -------------------------- & --------------- & -------------------- \\
多项式函数 & $a_n x^n + \cdots + a_0$ & $\mathbb{R}$ & 连续可导，曲线光滑 \\
有理函数 & $\frac{P(x)}{Q(x)}$ & $Q(x) \neq 0$ & 可能有渐近线和间断点 \\
指数函数 & $a^x \; (a > 0)$ & $\mathbb{R}$ & 单调，恒正 \\
对数函数 & $\log_a x$ & $x > 0$ & 指数函数的反函数 \\
三角函数 & $\sin x, \cos x$ & $\mathbb{R}$ & 周期函数，有界 \\
\bottomrule
\end{tabularx}
\end{table}

\section{极限理论}

\subsection{极限的定义}

\begin{quote}
\textbf{$\varepsilon$-$\delta$ 定义}：设函数 $f(x)$ 在点 $x_0$ 的某个去心邻域内有定义。若 $\forall \varepsilon > 0$，$\exists \delta > 0$，使得当 $0 < |x - x_0| < \delta$ 时，恒有 $|f(x) - A| < \varepsilon$，则称 $A$ 为 $f(x)$ 当 $x \to x_0$ 时的极限，记作：
\end{quote}

$$
\lim_{x \to x_0} f(x) = A
$$

\subsection{极限的四则运算法则}

设 $\lim f(x) = A$，$\lim g(x) = B$，则：

\begin{enumerate}
  \item $\lim [f(x) \pm g(x)] = A \pm B$
  \item $\lim [f(x) \cdot g(x)] = A \cdot B$
  \item $\lim \frac{f(x)}{g(x)} = \frac{A}{B} \quad (B \neq 0)$
\end{enumerate}

\subsection{重要极限公式}

两个必须牢记的重要极限：

\begin{enumerate}
  \item 第一重要极限：
\end{enumerate}

$$
\lim_{x \to 0} \frac{\sin x}{x} = 1
$$

\begin{enumerate}
  \item 第二重要极限：
\end{enumerate}

$$
\lim_{x \to \infty} \left(1 + \frac{1}{x}\right)^x = e
$$

\section{导数与微分}

\subsection{导数的定义}

设函数 $y = f(x)$ 在点 $x_0$ 的某个邻域内有定义，当自变量 $x$ 在 $x_0$ 处取得增量 $\Delta x$ 时，相应地函数取得增量 $\Delta y = f(x_0 + \Delta x) - f(x_0)$。如果极限

$$
f'(x_0) = \lim_{\Delta x \to 0} \frac{\Delta y}{\Delta x} = \lim_{\Delta x \to 0} \frac{f(x_0 + \Delta x) - f(x_0)}{\Delta x}
$$

存在，则称函数 $f(x)$ 在点 $x_0$ 处\textbf{可导}，该极限值称为 $f(x)$ 在 $x_0$ 处的\textbf{导数}。

\subsection{基本求导公式}

\begin{table}[htbp]
\centering
\begin{tabularx}{\textwidth}{>{\hsize=0.929\hsize\linewidth=\hsize}X>{\hsize=1.071\hsize\linewidth=\hsize}X>{\hsize=1.000\hsize\linewidth=\hsize}X}
\toprule
\textbf{函数$f(x)$} & \textbf{导数$f'(x)$} & \textbf{注释} \\
\midrule
------------- & --------------- & -------------- \\
$C$（常数） & $0$ & 常数变化率为零 \\
$x^n$ & $nx^{n-1}$ & 幂函数求导公式 \\
$\sin x$ & $\cos x$ & 正弦求余弦 \\
$\cos x$ & $-\sin x$ & 余弦求负正弦 \\
$\ln x$ & $\frac{1}{x}$ & 对数函数的导数 \\
\bottomrule
\end{tabularx}
\end{table}

\section{定积分}

\subsection{定积分的定义}

定积分来源于计算曲边梯形的面积。将区间 $[a, b]$ 分成 $n$ 个小区间：

$$
\int_a^b f(x)\,dx = \lim_{n \to \infty} \sum_{i=1}^{n} f(\xi_i) \Delta x_i
$$

其中 $\xi_i$ 是第 $i$ 个小区间内的任意点，$\Delta x_i$ 是第 $i$ 个小区间的长度。

\subsection{牛顿-莱布尼茨公式}

若 $F'(x) = f(x)$，则：

$$
\int_a^b f(x)\,dx = F(b) - F(a) = \big[F(x)\big]_a^b
$$

此公式将积分计算转化为寻找原函数的问题，是微积分中最重要的公式之一。

\subsection{常用积分表}

\begin{table}[htbp]
\centering
\begin{tabularx}{\textwidth}{>{\hsize=0.898\hsize\linewidth=\hsize}X>{\hsize=1.102\hsize\linewidth=\hsize}X}
\toprule
\textbf{被积函数$f(x)$} & \textbf{原函数$\int f(x)\,dx$} \\
\midrule
---------------------- & --------------------------- \\
$k$（常数） & $kx + C$ \\
$x^n \; (n \neq -1)$ & $\frac{x^{n+1}}{n+1} + C$ \\
$\frac{1}{x}$ & $\ln|x| + C$ \\
$e^x$ & $e^x + C$ \\
$a^x$ & $\frac{a^x}{\ln a} + C$ \\
\bottomrule
\end{tabularx}
\end{table}

\end{document}
